\documentclass[11pt]{article}
\usepackage{mystyle}

\title{Rosen --- \S 5.3: Recursive Definitions and Structural Induction\\ \large Selected Exercises}
\author{Alston}
\date{Semester 2, 2026}

\begin{document}
\maketitle

% ======================================================================
\begin{exercise}{14}
    \textit{In this and the next exercise, $f_n$ is the $n$th
    Fibonacci number.}

    Show that $f_{n+1} f_{n-1} - f_n^2 = (-1)^n$ when $n$ is a
    positive integer.
\end{exercise}

% ======================================================================
\begin{exercise}{18}
    Let
    \[
    A = \begin{bmatrix} 1 & 1 \\ 1 & 0 \end{bmatrix}.
    \]
    Show that
    \[
    A^n = \begin{bmatrix} f_{n+1} & f_n \\ f_n & f_{n-1} \end{bmatrix}
    \]
    when $n$ is a positive integer.
\end{exercise}

% ======================================================================
\begin{exercise}{20}
    Give a recursive definition of the functions $\max$ and $\min$ so
    that $\max(a_1, a_2, \dots, a_n)$ and $\min(a_1, a_2, \dots, a_n)$
    are the maximum and minimum of the $n$ numbers
    $a_1, a_2, \dots, a_n$, respectively.
\end{exercise}

% ======================================================================
\begin{exercise}{21}
    Let $a_1, a_2, \dots, a_n$, and $b_1, b_2, \dots, b_n$ be real
    numbers. Use the recursive definitions that you gave in Exercise
    20 to prove these.
    \begin{enumerate}[label=(\alph*)]
        \item $\max(-a_1, -a_2, \dots, -a_n) = -\min(a_1, a_2, \dots,
        a_n)$
        \item $\max(a_1+b_1, a_2+b_2, \dots, a_n+b_n) \leq
        \max(a_1, a_2, \dots, a_n) + \max(b_1, b_2, \dots, b_n)$
        \item $\min(a_1+b_1, a_2+b_2, \dots, a_n+b_n) \geq
        \min(a_1, a_2, \dots, a_n) + \min(b_1, b_2, \dots, b_n)$
    \end{enumerate}
\end{exercise}

% ======================================================================
\begin{exercise}{26}
    Let $S$ be the subset of the set of ordered pairs of integers
    defined recursively by

    \textbf{Basis step:} $(0,0) \in S$.

    \textbf{Recursive step:} If $(a,b) \in S$, then $(a+2, b+3) \in S$
    and $(a+3, b+2) \in S$.
    \begin{enumerate}[label=(\alph*)]
        \item List the elements of $S$ produced by the first five
        applications of the recursive definition.
        \item Use strong induction on the number of applications of
        the recursive step of the definition to show that
        $5 \mid a+b$ when $(a,b) \in S$.
        \item Use structural induction to show that $5 \mid a+b$ when
        $(a,b) \in S$.
    \end{enumerate}
\end{exercise}

% ======================================================================
\begin{exercise}{32}
    \begin{enumerate}[label=(\alph*)]
        \item Give a recursive definition of the function $ones(s)$,
        which counts the number of ones in a bit string $s$.
        \item Use structural induction to prove that
        $ones(st) = ones(s) + ones(t)$.
    \end{enumerate}
\end{exercise}

% ======================================================================
\begin{exercise}{36}
    \textit{The reversal of a string is the string consisting of the
    symbols of the string in reverse order. The reversal of the
    string $w$ is denoted by $w^R$.}

    Use structural induction to prove that
    $(w_1 w_2)^R = w_2^R w_1^R$.
\end{exercise}

% ======================================================================
\begin{exercise}{38}
    Give a recursive definition of the set of bit strings that are
    palindromes.
\end{exercise}

% ======================================================================
\begin{exercise}{44}
    \textit{The set of leaves and the set of internal vertices of a
    full binary tree can be defined recursively.}

    \textit{\textbf{Basis step:} The root $r$ is a leaf of the full
    binary tree with exactly one vertex $r$. This tree has no
    internal vertices.}

    \textit{\textbf{Recursive step:} The set of leaves of the tree
    $T = T_1 \cdot T_2$ is the union of the sets of leaves of $T_1$
    and of $T_2$. The internal vertices of $T$ are the root $r$ of
    $T$ and the union of the set of internal vertices of $T_1$ and
    the set of internal vertices of $T_2$.}

    Use structural induction to show that $l(T)$, the number of
    leaves of a full binary tree $T$, is 1 more than $i(T)$, the
    number of internal vertices of $T$.
\end{exercise}

% ======================================================================
\begin{exercise}{47}
    A \emph{partition} of a positive integer $n$ is a way to write
    $n$ as a sum of positive integers where the order of terms in
    the sum does not matter. For instance, $7 = 3+2+1+1$ is a
    partition of 7. Let $P_m$ equal the number of different
    partitions of $m$, and let $P_{m,n}$ be the number of different
    ways to express $m$ as the sum of positive integers not
    exceeding $n$.
    \begin{enumerate}[label=(\alph*)]
        \item Show that $P_{m,m} = P_m$.
        \item Show that the following recursive definition for
        $P_{m,n}$ is correct:
        \[
        P_{m,n} =
        \begin{cases}
            1 & \text{if } m = 1 \\
            1 & \text{if } n = 1 \\
            P_{m,m} & \text{if } m < n \\
            1 + P_{m,m-1} & \text{if } m = n > 1 \\
            P_{m,n-1} + P_{m-n,n} & \text{if } m > n > 1.
        \end{cases}
        \]
        \item Find the number of partitions of 5 and of 6 using this
        recursive definition.
    \end{enumerate}
\end{exercise}

% ======================================================================
\begin{exercise}{58}
    Show that each of these proposed recursive definitions of a
    function on the set of positive integers does not produce a
    well-defined function.
    \begin{enumerate}[label=(\alph*)]
        \item $F(n) = 1 + F(\lfloor n/2 \rfloor)$ for $n \geq 1$ and
        $F(1) = 1$.
        \item $F(n) = 1 + F(n-3)$ for $n \geq 2$, $F(1) = 2$, and
        $F(2) = 3$.
        \item $F(n) = 1 + F(n/2)$ for $n \geq 2$, $F(1) = 1$, and
        $F(2) = 2$.
        \item $F(n) = 1 + F(n/2)$ if $n$ is even and $n \geq 2$,
        $F(n) = 1 - F(n-1)$ if $n$ is odd, and $F(1) = 1$.
        \item $F(n) = 1 + F(n/2)$ if $n$ is even and $n \geq 2$,
        $F(n) = F(3n-1)$ if $n$ is odd and $n \geq 3$, and
        $F(1) = 1$.
    \end{enumerate}
\end{exercise}

% ======================================================================
\begin{exercise}{53}
    \textit{Recall Ackermann's function:}
    \[
    A(m,n) =
    \begin{cases}
        2n & \text{if } m = 0 \\
        0 & \text{if } m \geq 1 \text{ and } n = 0 \\
        2 & \text{if } m \geq 1 \text{ and } n = 1 \\
        A(m-1, A(m, n-1)) & \text{if } m \geq 1 \text{ and } n \geq 2.
    \end{cases}
    \]

    Prove that $A(m, n+1) > A(m, n)$ whenever $m$ and $n$ are
    nonnegative integers.
\end{exercise}

\end{document}